\section{Related Work}

Graph-based interpolation for missing data in structured signals is well established in the signal processing on graphs literature. Early formulations by Narang and Ortega \cite{narang2013} showed that interpolation can be cast as a graph-constrained reconstruction problem, while later work by Puy \emph{et al.} \cite{puy2018} clarified when bandlimited graph signals can be recovered from partial observations. In parallel, the broader GSP framing summarized by Shuman \emph{et al.} \cite{shuman2013} and later graph-wavelet constructions by Hammond \emph{et al.} \cite{hammond2011} helped establish the intuition that graph structure can replace purely Euclidean interpolation when the data lives on irregular domains.

More recent methods include RKHS-based recovery schemes such as BGSRP \cite{zhang2024bgsrp}, as well as time-varying formulations that combine spatial and temporal smoothness \cite{giraldo2022}. In EEG settings, these ideas are especially relevant because electrodes are spatially arranged, the geometry is not perfectly uniform, and physiological activity evolves over time in a correlated but not trivial way. For this reason, methods that exploit both graph smoothness and temporal regularity are often more appropriate than simple baselines when channels are missing.

The present work does not attempt to introduce a new interpolation theory. Instead, it tries to make the comparison fairer and more reproducible by evaluating these families under shared data, masking, and metric protocols, using the same EEG benchmark pipeline and the same consolidated evidence artifacts.
